and the exact @@M0@@ sequence obtained by tensoring with @@M1@@:
@@M2@@
The image of the morphism: @@M3@@ is
@@M4@@. Consequently, we obtain the exact sequence:
@@M5@@

Then tensor it with @@M6@@; since
@@M7@@ is an @@M8@@-Module, and
@@M9@@ is invertible, we obtain the exact sequence:
\sgaStableTarget{sga6.e07.eq.3-2-4}@@M10@@
One readily checks that the morphism @@M11@@ occurring in this sequence is the canonical morphism, so that this sequence is identified with (3.1.2). By definition, we therefore have:
@@M12@@

\sgaStableTarget{sga6.e07.para.3-3}\textbf{3.3.} We now suppose that @@M13@@ is a scheme having an ample sheaf. Then the same is true of @@M14@@, which is projective over @@M15@@, and of @@M16@@ and @@M17@@. Furthermore, @@M18@@ and @@M19@@ are perfect morphisms, since they are regular immersions; @@M20@@ is smooth and hence perfect, since @@M21@@; finally, @@M22@@ is perfect by 1.9. Since these morphisms are projective, we can define the additive homomorphisms (IV
@@M23@@
